En un ensayo de calentamiento se normalizan el tiempo, la temperatura en exceso y la potencia, de modo que las variables del modelo son adimensionales:
\[
x'=-x+u(t),\qquad x(0)=0.
\]
La entrada \(u\) es continua a trozos, \(0\le u(t)\le1\) en \([0,1]\), y se apaga para \(t>1\). El balance se interpreta por tramos, con \(x\) continua. El registrador sólo conserva el aporte total
\[
\int_0^1u(t)\,dt=\frac12,
\]
no su distribución temporal. Halle el intervalo más ajustado posible para la temperatura \(x(2)\). Justifique ambos extremos y muestre entradas admisibles que los alcancen. Explique por qué conocer el aporte total no basta para conocer la temperatura final.
